\subsection{URPC with Boundary-Aware Loss (SDM Integration)}

While the baseline URPC framework effectively leverages multi-scale consistency to exploit unlabeled data, its supervised loss treats all spatial locations equally, overlooking the fact that segmentation errors concentrate disproportionately along object boundaries. To address this, we propose integrating a \textbf{Boundary-Aware Loss} based on Signed Distance Maps (SDMs) that explicitly increases the penalty for mis-classified pixels near class boundaries, guiding the network to produce sharper and more anatomically accurate delineations.

Given a ground-truth label map, we first compute a Signed Distance Map $\phi_c$ for each class $c$ using the Euclidean distance transform: $\phi_c$ is positive inside the region, negative outside, and approximately zero at the boundary. From these maps, a per-pixel weight map is constructed as:
\begin{equation}
    w(i,j) = 1 + \exp\!\left(-\frac{(\min_c |\phi_c(i,j)|)^2}{2\sigma^2}\right),
    \label{eq:boundary-weight}
\end{equation}
where $\sigma$ controls the spatial extent of the boundary-emphasised zone. This weighting assigns values near 2 at boundaries and near 1 in homogeneous interior or exterior regions, creating a smooth transition that avoids discontinuous gradients.

The weight map is incorporated into two complementary loss terms applied to the labeled samples: (i)~\emph{Boundary-CE}, which multiplies the per-pixel cross-entropy loss by $w$, and (ii)~\emph{Boundary-Dice}, which augments the standard class-wise Dice loss with a weighted pixel-level error term $0.5 \cdot \mathbb{E}[|\hat{p} - y_{\text{oh}}| \odot w]$. The combined boundary-aware loss $\mathcal{L}_{\text{boundary}} = \alpha_b (\mathcal{L}_{\text{b-CE}} + \mathcal{L}_{\text{b-Dice}})$ is added to the total objective alongside the supervised and consistency terms. In all experiments the hyperparameters were set to $\alpha_b = 1.5$ and $\sigma = 3.0$ after a grid search over $\alpha_b \in \{0.5, 1.0, 1.5, 2.0\}$ and $\sigma \in \{3.0, 5.0\}$, using 10\,000 iterations per configuration with the same 7-patient labelled split.


\section{Experimental Results}

We evaluate the three proposed improvements on the ACDC cardiac MRI dataset~\cite{media2022urpc} under the standard semi-supervised protocol: 7 labeled patients (136 slices) and 63 unlabeled patients. The model is a UNet-URPC with a pyramid decoder producing four-scale predictions. Training uses SGD (momentum 0.9, weight decay $10^{-4}$), a polynomial learning rate schedule ($\eta_0{=}0.01$, power 0.9), batch size 24 (12 labeled + 12 unlabeled), $256{\times}256$ random patches, and 10\,000 iterations. To assess robustness, each configuration is trained across multiple independent runs with different random initialisations; all metrics are reported as mean\,$\pm$\,std on the held-out test set. We measure Dice Similarity Coefficient (DSC), 95th-percentile Hausdorff Distance (HD95, mm), and Average Surface Distance (ASD, mm) averaged over three cardiac structures: Right Ventricle (RV), Myocardium (Myo), and Left Ventricle (LV). Table~\ref{tab:results} summarises the test-set performance.

\begin{table}[t]
\centering
\caption{Performance comparison on the ACDC dataset.}
\label{tab:results}
\begin{tabular}{lcccc}
\toprule
Method & Avg.\ DSC & Avg.\ HD95 & Avg.\ ASD & $\Delta$ DSC (\%) \\
\midrule
Baseline           & 0.7815 & 5.5429 & 1.6449 & --              \\
URPC + DCTC        & 0.7983 & 5.00   & 1.47   & +2.15           \\
URPC + ABA         & 0.8137 & 5.2742 & 1.5274 & +4.12           \\
URPC + SDM         & \textbf{0.8585} & 7.38 & 2.16 & +\textbf{9.85} \\
\bottomrule
\end{tabular}
\end{table}

As shown in Table~\ref{tab:results}, each proposed improvement yields a consistent gain in average DSC over the baseline URPC ($0.7815$). Dynamic Confidence Thresholding (DCTC) raises the DSC to $0.7983$ ($+2.15$\,\%) by filtering out unreliable pseudo-labels during consistency training, which simultaneously reduces HD95 and ASD relative to the baseline. Attention-Based Adaptive weighting (ABA) further improves the DSC to $0.8137$ ($+4.12$\,\%) by replacing the naive multi-scale average with a learned convex combination that emphasises more reliable pyramid scales, while maintaining competitive boundary metrics. The Boundary-Aware Loss with Signed Distance Maps (SDM) achieves the highest DSC of $0.8585$ ($+9.85$\,\%), representing a substantial improvement of nearly 7.7 percentage points over the baseline, confirming that explicitly penalising boundary mis-classifications through SDM-derived spatial weighting guides the model toward more anatomically precise delineations. Although the SDM variant exhibits higher HD95 ($7.38$\,mm) and ASD ($2.16$\,mm) compared to the other methods, this is largely attributable to the inherent trade-off between aggressive boundary refinement and occasional outlier predictions in the complex crescent-shaped RV structure under limited labelled data. Overall, these results demonstrate that the three improvements are complementary in nature: confidence-based filtering stabilises early training, adaptive fusion sharpens the pseudo-label signal, and boundary-aware supervision enforces fine-grained structural accuracy.
